% AI-written, AI cross-reviewed; no human sign-off; not a claim that RH is proved
% Atlas unit D568. Local manuscript v3, item G.602, heading on printed PDF p. 307.
% Formalism.LFunctionChainParameters.WeightedSlopeBounds.Ga_bound5
% paper2/decomposition_tables.tex:4992–4997. Verbatim excerpt; SHA-256 feff4980b86a3fcdb87284017a9edeeb7af62c53183f91a86071571b1c9c2431.
% Original TeX uses the macros and cross-reference keys of the source manuscript.

\begin{formalstatement}[D568]\label{decomp:D568}
Under the ten size, division, comparison, counting, slope, height and finite-parameter assumptions, Ga(u) is at most Main(u)+log|W(u)|+K0v5+M plus (1\ensuremath{-}delBar) inverse times [4 log(1+x\textasciicircum{}cHH)(K3+M+b)+(etaBar+4sqrt(eBar))x\textasciicircum{}cHH b], where M is the weighted mean of Lo5.
\end{formalstatement}
{\small\raggedright\noindent\textsf{\bfseries Project-specific}\quad \nolinkurl{Formalism.LFunctionChainParameters.WeightedSlopeBounds.Ga_bound5}\par\noindent \nolinkurl{Formalism/SizeWindowBrackets/ChainSlopeWindows/LevelOfCongruence.lean:113--160}\par\noindent SHA-256 \nolinkurl{17c66172050c9fc07489b18ab7c496528895fc4d448e9c288f8cfb0c1ce40f0a}\par\noindent Dominated 233\quad Disjoint 233\quad terminal.\par}
Arithmetic bundle slope estimates do not compare the specified Ga, weighted mean and Main with these explicit loss terms.
\hyperref[lit:S064]{Published sources and search S064}.
